% Generated draft: requires mathematical and source review.
\documentclass[11pt]{article}
\usepackage[margin=25mm]{geometry}
\usepackage{fontspec}
\setmainfont{Latin Modern Roman}
\usepackage{amsmath,amssymb,mathtools,unicode-math}
\setmathfont{Latin Modern Math}
\usepackage{enumitem,needspace,xurl,hyperref}
\setlength{\parindent}{0pt}
\setlength{\parskip}{5pt}
\setlength{\emergencystretch}{3em}
\newcommand{\N}{\mathbb{N}}
\newcommand{\Z}{\mathbb{Z}}
\newcommand{\Q}{\mathbb{Q}}
\newcommand{\R}{\mathbb{R}}
\newcommand{\C}{\mathbb{C}}
\newcommand{\F}{\mathbb{F}}
\newcommand{\A}{\mathbb{A}}
\newcommand{\PP}{\mathbb{P}}
\newcommand{\E}{\mathbb{E}}
\newcommand{\eps}{\varepsilon}
\newcommand{\dd}{\,\mathrm{d}}
\begin{document}
\section*{UnsolvedMath ID: 30006064}
\textbf{Source problem code:} catalogue entry
\subsection{Definitions and mathematical statement}
Let \(A=\bigoplus_{i\geq 0}A_i\) be an Artin--Schelter regular graded algebra, as stipulated by the supplied reconstructed statement, and let
\[
Z(A)=\{z\in A:za=az\text{ for every }a\in A\}.
\]
For a homogeneous element \(f\in Z(A)\cap A_d\), put \(B=A/(f)\), where \((f)\) is the two-sided ideal generated by \(f\). The record does not specify the base field, the precise convention for Artin--Schelter regularity, the dimension of \(A\), or whether \(d>0\), \(f\neq 0\), and multiplication by \(f\) is injective. These are unresolved hypotheses; they are not imposed here without qualification.

Under the usual locally finite, connected graded, right noetherian hypotheses over a field \(k\), one defines
\[
\operatorname{qgr}A=\operatorname{grmod}A/\operatorname{fdim}A,
\]
where \(\operatorname{grmod}A\) is the category of finitely generated graded right \(A\)-modules with degree-preserving morphisms, and \(\operatorname{fdim}A\) is its Serre subcategory of modules finite dimensional over \(k\). Here connected means \(A_0=k\), and locally finite means \(\dim_k A_i<\infty\) for every \(i\). The analogous definition applies to \(B\). This explains a standard interpretation of the notation, rather than asserting that the supplied record states all these hypotheses. With this interpretation, restriction of scalars along \(A\to B\) gives the hypersurface category denoted in the source by
\[
\operatorname{qgr}(A/(f))\subseteq\operatorname{qgr}A.
\]

\textbf{Problem.} For the intended class of Artin--Schelter regular graded algebras \(A\) and homogeneous central elements \(f\), classify the noncommutative hypersurfaces \(\operatorname{qgr}(A/(f))\), describe their singularities, and describe their derived categories, extending the understanding available in quadratic and finite-over-the-center cases. This is a classification programme, not a conjecture asserting a specified equivalence or invariant. The record leaves the classification equivalence relation, the meaning of singularity or smoothness, and the choice of derived category unspecified. For example, \(D^b(\operatorname{qgr}B)\) denotes the bounded derived category of the abelian category \(\operatorname{qgr}B\), but the record does not explicitly select this construction.

The supplied excerpt reports extensive work for quadrics, meaning \(d=2\): a Koszul duality criterion for smoothness in its reference [1], descriptions of derived categories in some examples in [2], and a full classification for quadrics in noncommutative surfaces, denoted \(n=2\), in [3]. The meaning of \(n\) and the hypotheses of these results are absent. These are supplied background claims, not assertions verified by this draft.

For \(A\) finite over \(Z(A)\), the excerpt attributes to its reference [5] a construction involving a variety \(Y\), a finite morphism \(Y\to\operatorname{Proj}Z(A)\), and a sheaf of maximal orders \(\mathcal A\), with reported equivalences
\[
\operatorname{coh}(Y,\mathcal A)\simeq\operatorname{qgr}A,
\qquad
\operatorname{coh}(V,\mathcal A|_V)\simeq\operatorname{qgr}(A/(f)).
\]
Here \(\operatorname{coh}\) denotes coherent right modules over the indicated sheaf of algebras, and \(V\subseteq Y\) is the reported vanishing locus associated with \(f\). The excerpt describes \(\mathcal A\) as Azumaya on a dense locus, with ramification locus \(\Delta\subseteq Y\), and describes restrictions away from \(\Delta\) as behaving like possibly twisted ordinary hypersurfaces. The hypotheses needed for these statements and equivalences are not supplied. In particular, the excerpt calls the object defined by homogeneous \(f\) a function on \(Y\); its projective interpretation may require a section of a twist. The excerpt ends mid-sentence, so no further condition on \(V\) can be recovered.
\subsection{Short English statement}
The supplied problem asks for a classification of noncommutative projective hypersurfaces obtained by quotienting an Artin--Schelter regular graded algebra by a homogeneous central element, together with descriptions of their singularities and derived categories. It seeks progress beyond studied quadratic and finite-over-the-center settings. Several hypotheses and the intended classification criteria are missing, and the original excerpt is truncated. The supplied open status is dated curation, not an independently verified conclusion.
\subsection{Sources}
\begin{itemize}
\item Supplied local UnsolvedMath registry record, permanent ID \(30006064\), title ``Classification of Noncommutative Hypersurfaces,'' source problem code \texttt{catalogue entry}, in the Oberwolfach Reports Open Problems collection. No proposer is identified. The record gives proposed year 2024.
\item Supplied source citation: \emph{MFO-RIMS Tandem Workshop: Algebraic Geometry and Noncommutative Projective Varieties (2025)}, Oberwolfach Reports, pp.\ 2551--2606, DOI \texttt{10.4171/owr/2024/44}. Supplied source and literature link, labelled ``Oberwolfach Report 44/2024'': \url{https://doi.org/10.4171/owr/2024/44}. The citation's 2025 workshop label differs from the report's 2024 designation; this discrepancy is unresolved.
\item The supplied original excerpt uses references [1], [2], [3], and [5], but does not provide their bibliographic entries or identify the speaker whose paper is described as [5]. Their identities and precise results cannot be reconstructed from this record.
\item Literature assessment dated exactly \texttt{2026-08-22}: the supplied curation labels the problem open and reports important special-case results for quadrics and finite-over-the-center orders, while describing the general classification and derived-category programme as open. This is supplied dated curation, not independently verified. The background explicitly requests independent verification of time-sensitive status and cited claims.
\item The registry labels its statement \texttt{corrected\_verified} and says it was reconstructed and checked against the official report. These are supplied provenance claims; this draft does not independently verify them or claim to have visited the report.
\end{itemize}
\end{document}
